\clearpage
\section{Full term weakening and arbitrary finite environments}
\label{sec:term-semantic-weakening}
The semantic weakening induction is now complete for the original, full
intrinsically scoped term language. For every term and either kind of unused
resource, its actual renamed translation is connected to an outside weakening
of its original translation by a finite calculation of proved Idem and BPullW
fillings and recursive graph isomorphisms. The two unchanged Proposition~12.2
programs consequently satisfy this law in every finite pair of environments,
including the named distinct-list boundary. These are proved calculations,
not hypotheses asserting that translation respects weakening.

The milestone finishes the weakening task in Appendix~B.1. It does not prove
the equivalence between the two closed programs: Figure~22's intermediate
graphs and their equivalence chain still have to be assembled, along with the
remaining correspondence obligations. The original terms, operation alphabet,
translation clauses and machine rules are unchanged.

\paragraph{Statement and notation.}
Let $E_{k,h}=*^k\diamond^h$ and write $W_sG$ for adjoining one primitive
weakening at the boundary between the term and atom blocks. For $s=*$ the
variable injection is $i\mapsto i$ into $\mathrm{Fin}(k+1)$ and atoms are
unchanged. For $s=\diamond$ variables are unchanged and the atom injection is
$a\mapsto a+1$ into $\mathrm{Fin}(h+1)$. Write these maps as $\rho_s,\alpha_s$.
The theorem constructs
\[
 \lbrack\!\lbrack\operatorname{rename}(\rho_s,\alpha_s,t)\rbrack\!\rbrack
 \leadsto W_s\lbrack\!\lbrack t\rbrack\!\rbrack.
\]
Here the left graph's depth index is transported along the already proved
renaming-depth equality. This changes only an equality index, not the graph.
The calculation symbol has its existing meaning: isomorphisms, actual Idem
and BPullW fillings, symmetry and transitivity. Its semantic interpretation
preserves exact successful-run counts at the admitted recursive closing sites.
No evaluation or termination premise is imposed on $t$.

\Needspace{12\baselineskip}
\subsection{Occurrences under ordered graph constructions}

Equal-interface transport of an occurrence is defined by eliminating the
four equalities identifying its source and target boundaries. The input and
output laws transport the corresponding finite indices by the same equalities;
after equality elimination, both proofs are the supplied port equation.

\lean{Copying/OccurrenceInterfaceTransport}{7}{11}{Opening.Occurrence.castInterfaces}

\lean{Copying/OccurrenceInterfaceTransport}{12}{18}{Opening.Occurrence.castInterfaces_input}

\lean{Copying/OccurrenceInterfaceTransport}{19}{24}{Opening.Occurrence.castInterfaces_output}

The occurrence of a tensor component is obtained from the actual tagged
parallel inclusion. Two occurrences combine by the parallel-join constructor:
their target edge images have different Boolean tags, so are disjoint. The
left and right port laws follow from the join's component equation. Its four
boundary specializations use the original tensor joining map, so both input
and output orders are retained rather than inferred from their cardinalities.

\lean{Copying/OccurrenceTensorMap}{8}{11}{Opening.Occurrence.tensorComponent}

\lean{Copying/OccurrenceTensorMap}{12}{24}{Opening.Occurrence.tensorMap}

\lean{Copying/OccurrenceTensorMap}{25}{33}{Opening.Occurrence.tensorMap_left}

\lean{Copying/OccurrenceTensorMap}{34}{42}{Opening.Occurrence.tensorMap_right}

\lean{Copying/OccurrenceTensorMap}{43}{52}{Opening.Occurrence.tensorMap_input_left}

\lean{Copying/OccurrenceTensorMap}{53}{62}{Opening.Occurrence.tensorMap_input_right}

\lean{Copying/OccurrenceTensorMap}{63}{72}{Opening.Occurrence.tensorMap_output_left}

\lean{Copying/OccurrenceTensorMap}{73}{83}{Opening.Occurrence.tensorMap_output_right}

A suffix occurrence fixing its input boundary can also be prefixed by an
unchanged graph. The serial join glues each output of that prefix to the same
suffix input; the fixed-input hypothesis is exactly the compatibility needed
at this shared boundary. Its input law fixes the exterior prefix inputs and
its output law follows the suffix's exterior output map.

\lean{Copying/OccurrenceTensorMap}{84}{93}{Opening.Occurrence.serialPrefix}

\lean{Copying/OccurrenceTensorMap}{94}{101}{Opening.Occurrence.serialPrefix_input}

\lean{Copying/OccurrenceTensorMap}{102}{110}{Opening.Occurrence.serialPrefix_output}

\Needspace{12\baselineskip}
\subsection{Injective maps of the literal atomic graphs}

In selected-resource coordinates there is one internal port, one input and
one output per resource. The port map fixes the first two sorts and applies
an arbitrary injection $\rho$ to outputs. The source-incidence equation is a
case split on the edge tag. The target equation additionally uses
$\rho(j)=\rho(i)\iff j=i$ to preserve the distinction between the selected row
and all other rows. Thus the actual selected contraction and its attached
weakening are retained; they are not replaced by an ideal identity wire.

\lean{Copying/SelectedResourceEmbedding}{10}{14}{SelectedResource.mapOutput}

\lean{Copying/SelectedResourceEmbedding}{15}{17}{SelectedResource.source_map}

\lean{Copying/SelectedResourceEmbedding}{18}{29}{SelectedResource.target_map}

The raw occurrence conjugates this coordinate map by the existing port
bijections. Internals use the one-internal-vertex bijections; edges use the
canonical index $\mathrm{Fin}(n)+1$ and the map $\rho+\mathrm{id}$. Injectivity
follows by disjoint-sum cases. Label preservation follows from the canonical
edge-label equations, and ordered incidences follow from the two preceding
map equations. The input-interface equality removes all zero-use blocks;
the output-interface equality sums the before, selected and after block
lengths to $n$. Transport along these equalities gives the occurrence between
the original selected-resource graphs. Its input is fixed and output $j$
is sent to output $\rho(j)$.

\lean{Copying/SelectedResourceEmbedding}{30}{65}{SelectedResource.rawEmbedding}

\lean{Copying/SelectedResourceEmbedding}{66}{69}{SelectedResource.input_interface}

\lean{Copying/SelectedResourceEmbedding}{70}{78}{SelectedResource.output_interface}

\lean{Copying/SelectedResourceEmbedding}{79}{84}{SelectedResource.embedding}

\lean{Copying/SelectedResourceEmbedding}{85}{93}{SelectedResource.embedding_input}

\lean{Copying/SelectedResourceEmbedding}{94}{103}{SelectedResource.embedding_output}

A zero-use distributor forest embeds under any injection by mapping its
weakening row $j$ to row $\rho(j)$. Each row has its original label and singleton
target, proving the output law. The variable occurrence tensors the selected
term-resource occurrence with the atom forest occurrence. The atom occurrence
retains its original instance edge in serial with the selected atom-resource
occurrence, alongside the term forest. The constant occurrence retains the
literal constant edge between the two forests. No semantic premise is used.

\lean{Copying/AtomicResourceEmbedding}{10}{39}{Weakening.Atomic.forestEmbedding}

\lean{Copying/AtomicResourceEmbedding}{40}{44}{Weakening.Atomic.forestEmbedding_output}

\lean{Copying/AtomicResourceEmbedding}{45}{52}{Weakening.Atomic.variableMap}

\lean{Copying/AtomicResourceEmbedding}{53}{61}{Weakening.Atomic.atom}

\lean{Copying/AtomicResourceEmbedding}{62}{69}{Weakening.Atomic.constant}

For each of these three atomic clauses, three separate boundary equations
are proved: the sole term input is fixed, term output $i$ goes to $\rho(i)$,
and atom output $a$ goes to $\alpha(a)$. The proofs apply the tensor and
serial port equations and then the exact original interface casts. Finite
indices are compared by their values, accounting explicitly for the length
of the preceding term block in atom coordinates.

\lean{Copying/AtomicEmbeddingBoundary}{10}{18}{Weakening.Atomic.variable_input}

\lean{Copying/AtomicEmbeddingBoundary}{19}{27}{Weakening.Atomic.variable_term}

\lean{Copying/AtomicEmbeddingBoundary}{28}{39}{Weakening.Atomic.variable_atom}

\lean{Copying/AtomicEmbeddingBoundary}{40}{52}{Weakening.Atomic.atom_input}

\lean{Copying/AtomicEmbeddingBoundary}{53}{62}{Weakening.Atomic.atom_term}

\lean{Copying/AtomicEmbeddingBoundary}{63}{77}{Weakening.Atomic.atom_atom}

\lean{Copying/AtomicEmbeddingBoundary}{78}{93}{Weakening.Atomic.constant_input}

\lean{Copying/AtomicEmbeddingBoundary}{94}{106}{Weakening.Atomic.constant_term}

\lean{Copying/AtomicEmbeddingBoundary}{107}{121}{Weakening.Atomic.constant_atom}

\Needspace{12\baselineskip}
\subsection{A complete one-edge complement}

\termdef{resource-gap}
A resource gap consists of an injective old-to-new index map, a fresh target
index outside its image, and a proof that every target index is either old or
that fresh index. Appending a variable uses $\mathrm{Fin.castSucc}$ with fresh
index $n$; prepending an atom uses successor with fresh index zero. The first
coverage proof splits a finite index at $n$, and the second splits it at zero.
Injectivity and exclusion of the fresh index follow from these index values.

\lean{Copying/ResourceGap}{6}{14}{Weakening.ResourceGap}

\lean{Copying/ResourceGap}{15}{24}{Weakening.ResourceGap.append}

\lean{Copying/ResourceGap}{25}{34}{Weakening.ResourceGap.prepend}

A literal weakening edge in a graph gives an occurrence of the primitive
weakening provided its label, empty source and singleton exterior target are
specified. The construction maps the unique edge to that edge and the unique
output to that target; its internal carrier is empty. The edge and output
lemmas record these defining equations.

\lean{Copying/WeakeningEdgeOccurrence}{8}{24}{Weakening.weakeningAt}

\lean{Copying/WeakeningEdgeOccurrence}{25}{30}{Weakening.weakeningAt_edge}

\lean{Copying/WeakeningEdgeOccurrence}{31}{34}{Weakening.weakeningAt_output}

\termdef{one-weakening-complement}
A complement certificate partitions the target's edges into a complete-label
occurrence of the old graph and an occurrence of one genuine primitive
weakening. Their images are disjoint and jointly exhaustive. The boundaries
may initially differ; ordered boundary conditions are supplied only when
constructing the final isomorphism. This is geometric evidence, not an assumed
weakening or observational-equivalence law.

\lean{Copying/WeakeningComplement}{10}{18}{Weakening.Complement}

Tensoring this certificate on either side retains the unrelated component
entirely in the old occurrence. Edge-tag cases prove disjointness and coverage:
an edge in the retained component is old, while an edge in the changed
component uses the original partition. Four interface equalities transport
the certificate by equality elimination. The old-occurrence equation and
added-output equation ensure the same port maps survive this transport.

\lean{Copying/WeakeningComplement}{19}{35}{Weakening.Complement.tensorRight}

\lean{Copying/WeakeningComplement}{36}{52}{Weakening.Complement.tensorLeft}

\lean{Copying/WeakeningComplement}{53}{57}{Weakening.Complement.castInterfaces}

\lean{Copying/WeakeningComplement}{58}{61}{Weakening.Complement.castInterfaces_old}

\lean{Copying/WeakeningComplement}{62}{68}{Weakening.Complement.castInterfaces_added_output}

Prefixing by a fixed serial graph retains all prefix edges in the old part.
The suffix's old input map must be fixed, so the serial gluing agrees. Its
coverage and disjointness proofs again split the serial edge tags. The added
weakening's output equations for right tensor, left tensor and serial prefix
are obtained by composing its given output equation with the corresponding
actual component insertion.

\lean{Copying/WeakeningComplement}{69}{87}{Weakening.Complement.serialPrefix}

\lean{Copying/WeakeningComplement}{104}{111}{Weakening.Complement.tensorRight_added_output}

\lean{Copying/WeakeningComplement}{112}{119}{Weakening.Complement.tensorLeft_added_output}

\lean{Copying/WeakeningComplement}{120}{127}{Weakening.Complement.serialPrefix_added_output}

A zero-use forest with a resource gap has precisely this partition: the
old rows are the injected rows and the sole remaining row is the weakening
at the gap. For the selected-resource graph, the gap cannot equal the selected
index, by exclusion from the image. Its edge therefore has the weakening
label and the fresh exterior target. Canonical edge-index case analysis proves
coverage: every ordinary row is old or fresh, and the distinguished selected
edge is always old. Equal-interface transport gives the certificate for the
original selected-resource clause, with its old occurrence equal to the
previously constructed embedding and its added output at the fresh index.

\lean{Copying/AtomicResourceComplement}{8}{26}{Weakening.Atomic.forestComplement}

\lean{Copying/AtomicResourceComplement}{27}{32}{Weakening.Atomic.forestComplement_output}

\lean{Copying/AtomicResourceComplement}{33}{71}{Weakening.Atomic.selectedRawComplement}

\lean{Copying/AtomicResourceComplement}{72}{77}{Weakening.Atomic.selectedComplement}

\lean{Copying/AtomicResourceComplement}{78}{81}{Weakening.Atomic.selectedComplement_old}

\lean{Copying/AtomicResourceComplement}{82}{91}{Weakening.Atomic.selectedComplement_output}

\Needspace{12\baselineskip}
\subsection{From the complete partition to a boundary-fixing isomorphism}

\termdef{weakening-extension-witness}
An extension witness adds three ordered boundary requirements to the
complement: old inputs are fixed, old outputs use the prescribed insertion
map, and the added weakening ends at the new position $k$. Its two components
are the old and added occurrences. Their pairwise edge disjointness follows
from the certificate, reversing the inequality for the opposite tag order.
The ordered parallel join and original interface transport then give an
occurrence of $W_sG$ in the target.

\lean{Copying/WeakeningExtensionIso}{10}{22}{Weakening.Extension.Witness}

\lean{Copying/WeakeningExtensionIso}{23}{25}{Weakening.Extension.components}

\lean{Copying/WeakeningExtensionIso}{26}{33}{Weakening.Extension.disjoint}

\lean{Copying/WeakeningExtensionIso}{34}{38}{Weakening.Extension.occurrence}

The old-port equation follows from the parallel join's component equation;
the input law follows by applying it to the retained payload input. For output
ports, first prove the equation separately for the old and added component.
The old case uses the insertion-index equation; the added case uses uniqueness
of the element of $\mathrm{Fin}(1)$ and the fact that the new position is $k$.
Every exterior output is recovered using the inverse of the ordered boundary
bijection, so the assembled occurrence fixes it pointwise. Edge surjectivity
is exactly the certificate's coverage, choosing the corresponding component
tag. The existing finite-occurrence theorem therefore yields a recursive
isomorphism. Its hypotheses include the exterior boundary equations just
proved, not merely an edge-count equality.

\lean{Copying/WeakeningExtensionIso}{39}{43}{Weakening.Extension.old_port}

\lean{Copying/WeakeningExtensionIso}{44}{46}{Weakening.Extension.input}

\lean{Copying/WeakeningExtensionIso}{47}{65}{Weakening.Extension.component_output}

\lean{Copying/WeakeningExtensionIso}{66}{75}{Weakening.Extension.output}

\lean{Copying/WeakeningExtensionIso}{76}{81}{Weakening.Extension.edge_surjective}

\lean{Copying/WeakeningExtensionIso}{82}{84}{Weakening.Extension.iso}

The complement-to-extension function packages those three proved boundary
equations into an extension witness and invokes its isomorphism. All atomic
applications below construct the complement and discharge the equations;
none assumes the desired isomorphism as an input.

\lean{Copying/WeakeningComplement}{88}{103}{Weakening.Complement.extensionIso}

\Needspace{12\baselineskip}
\subsection{Atomic base cases with the original selected rows}

Several identity equations identify retained components without replacing
their finite carriers. Occurrences with the same port and edge maps are equal:
the internal-map law determines the internal map, and remaining proof fields
are propositionally irrelevant. Tensoring two identity occurrences is the
identity by component port equations and edge tags; the serial-prefix identity
also checks the shared ports against the prefix's output embedding. The
forest embedding at identity is pointwise identity. For the raw selected
embedding, the port, internal and edge-coordinate inverse equations give the
same conclusion. Equality elimination proves the interface-cast identity,
and hence the selected-resource identity.

\lean{Copying/AtomicEmbeddingIdentity}{9}{17}{Weakening.Atomic.occurrence_ext}

\lean{Copying/AtomicEmbeddingIdentity}{18}{48}{Weakening.Atomic.tensorMap_identity}

\lean{Copying/AtomicEmbeddingIdentity}{49}{64}{Weakening.Atomic.serialPrefix_identity}

\lean{Copying/AtomicEmbeddingIdentity}{65}{74}{Weakening.Atomic.forestEmbedding_id}

\lean{Copying/AtomicEmbeddingIdentity}{75}{92}{Weakening.Atomic.rawEmbedding_id}

\lean{Copying/AtomicEmbeddingIdentity}{93}{97}{Weakening.Atomic.castIdentity}

\lean{Copying/AtomicEmbeddingIdentity}{98}{103}{Weakening.Atomic.selectedEmbedding_id}

The four resource-map equations identify variable insertion with
$\mathrm{Fin.castSucc}$, atom insertion with successor, and both unaffected
maps with identity. For a variable clause, variable insertion uses its
selected-term complement while atom insertion uses the atom-forest complement.
For an atom clause, variable insertion uses the term forest; atom insertion
uses the selected-atom complement prefixed by the unchanged instance edge.
For a constant clause the changed forest supplies the complement and the
constant edge remains in the old part. In every case tensoring retains all
other components. The three old-occurrence equalities use the preceding
identity laws to identify these composite witnesses with the original
atomic embedding maps.

\lean{Copying/AtomicWeakeningComplements}{10}{10}{Weakening.Atomic.termResource_term}

\lean{Copying/AtomicWeakeningComplements}{11}{11}{Weakening.Atomic.termResource_atom}

\lean{Copying/AtomicWeakeningComplements}{12}{12}{Weakening.Atomic.atomResource_term}

\lean{Copying/AtomicWeakeningComplements}{13}{14}{Weakening.Atomic.atomResource_atom}

\lean{Copying/AtomicWeakeningComplements}{15}{25}{Weakening.Atomic.variableComplement}

\lean{Copying/AtomicWeakeningComplements}{26}{39}{Weakening.Atomic.atomComplement}

\lean{Copying/AtomicWeakeningComplements}{40}{52}{Weakening.Atomic.constantComplement}

\lean{Copying/AtomicWeakeningComplements}{53}{60}{Weakening.Atomic.variableComplement_old}

\lean{Copying/AtomicWeakeningComplements}{61}{68}{Weakening.Atomic.atomComplement_old}

\lean{Copying/AtomicWeakeningComplements}{69}{74}{Weakening.Atomic.constantComplement_old}

The three added-output equations follow the selected complement's fresh
output through serial, tensor and interface transports. Their global value is
$k$ for either sort. The environment-output lemma splits an arbitrary exterior
index at $k$: below it the term-boundary law applies, and above it the atom
law applies after subtracting $k$. The insertion-index lemmas then identify
this with the required global old-output map.

\lean{Copying/AtomicWeakeningOutputs}{10}{26}{Weakening.Atomic.variableComplement_output}

\lean{Copying/AtomicWeakeningOutputs}{27}{47}{Weakening.Atomic.atomComplement_output}

\lean{Copying/AtomicWeakeningOutputs}{48}{70}{Weakening.Atomic.constantComplement_output}

\lean{Copying/AtomicWeakeningOutputs}{71}{82}{Weakening.Atomic.environment_output}

The variable, atom and constant isomorphisms now follow by applying the
complement-to-extension theorem to their actual complements. Input, old-output
and fresh-output conditions are the proved equations above. Reversing each
isomorphism and supplying the already proved focus freedom of the original
atomic translation gives the three base calculations, at every box allowance
and for either resource sort. This is the missing atomic step in the full
syntactic induction.

\lean{Copying/AtomicWeakening}{10}{18}{Weakening.variable_iso}

\lean{Copying/AtomicWeakening}{19}{27}{Weakening.atom_iso}

\lean{Copying/AtomicWeakening}{28}{37}{Weakening.constant_iso}

\lean{Copying/AtomicWeakening}{38}{43}{Weakening.variable_extend}

\lean{Copying/AtomicWeakening}{44}{49}{Weakening.atom_extend}

\lean{Copying/AtomicWeakening}{50}{53}{Weakening.constant_extend}

\Needspace{12\baselineskip}
\subsection{Binder transport and the full structural induction}

The lifted variable inclusion is again the inclusion: a finite-index split
checks the new bound variable zero and each successor separately. Expanding
syntactic renaming therefore identifies renaming an abstraction with
abstracting the correspondingly renamed body, for both variable and atom
extension; the atom case uses the identity-lifting law.

\lean{Copying/WeakeningTermTransport}{9}{12}{Weakening.lift_castSucc}

\lean{Copying/WeakeningTermTransport}{13}{16}{Weakening.rename_lam_term}

\lean{Copying/WeakeningTermTransport}{17}{20}{Weakening.rename_lam_atom}

Depth transport commutes with unary translation and with abstraction,
using the successor equality for the latter. Both proofs eliminate the depth
equality. Equal terms have equal depth-transported translations: eliminate
the term equality, then use proof irrelevance of the depth equality. Combining
this fact with the two binder-renaming equations and abstraction transport
gives the two explicit abstraction-translation equations. This avoids
rewriting a term underneath an equality whose type itself depends on that
term. No semantic assumption is hidden in these dependent transports.

\lean{Copying/WeakeningTermTransport}{21}{25}{Weakening.castDepth_unary}

\lean{Copying/WeakeningTermTransport}{26}{31}{Weakening.castDepth_abstraction}

\lean{Copying/WeakeningTermTransport}{32}{37}{Weakening.castTerm_eq}

\lean{Copying/WeakeningTermTransport}{38}{46}{Weakening.castTerm_lam_term}

\lean{Copying/WeakeningTermTransport}{47}{55}{Weakening.castTerm_lam_atom}

For application and binary operations, independent equalities for the two
operand depths induce an equality of their maxima. Transporting the original
maximum-padded translation is equal to translating the separately transported
operands with their new maximum padding. Eliminating the two depth equalities
reduces both statements to reflexivity; the original padding construction and
its depth inequalities are retained.

\lean{Copying/WeakeningTermTransport}{56}{65}{Weakening.castDepth_application}

\lean{Copying/WeakeningTermTransport}{66}{74}{Weakening.castDepth_binary}

\termdef{full-term-weakening}
The full theorem is proved by structural induction on $t$, generalizing over
the resource sort. Variables, atoms and constants use the three actual base
calculations. A unary constructor transports depth and invokes the already
proved unary induction step. Application and every binary operation combine
the independent operand hypotheses by the proved padded steps; the maximum
depth transport equations identify their sources with the original translated
renamed term. In the abstraction case, split the resource sort, use the
corresponding binder-translation equation, and apply the proved abstraction
step to the body's induction hypothesis. The body's boundary transport is
reflexive after the sort is fixed. Thus every constructor is discharged,
including arbitrary nested binders, shared variables and all store operations.

\lean{Copying/TermWeakening}{12}{50}{Weakening.term_weakening}

\Needspace{12\baselineskip}
\subsection{The original programs in all finite environments}

A calculation survives adding the same unused resource to both endpoints:
parallel composition retains the actual primitive weakening and its focus
freedom, then the original insertion-interface equality restores the desired
boundary. This lifts the actual proof leaves and introduces no new execution
step. Iterating term extension from an empty variable block defines
$V_k(G)$, and iterating atom extension from an empty atom block defines $A_h(G)$.
The ambient extension of a closed graph is $A_h(V_k(G))$; its ordered interface
is exactly $E_{k,h}$.

\lean{Copying/PropositionEnvironmentWeakening}{11}{18}{Weakening.Calculation.extend}

\lean{Copying/PropositionEnvironmentWeakening}{19}{24}{Weakening.addVariables}

\lean{Copying/PropositionEnvironmentWeakening}{25}{30}{Weakening.addAtoms}

\lean{Copying/PropositionEnvironmentWeakening}{31}{33}{Weakening.ambient}

Write $P_{k,h}$ and $Q_{k,h}$ for the unchanged allocating and constant-result
translations in Proposition~12.2. Applying full term weakening directly gives
$P_{k',h'}\leadsto W_sP_{k,h}$ and $Q_{k',h'}\leadsto W_sQ_{k,h}$.
These specializations need no extra term equivalence: the closed program
syntax is unchanged by ambient renaming, and its depth is respectively three
and two. Induction on $k$ proves the variable-only consequences. The base
calculation is identity; the successor composes the one-step theorem with
the previous calculation extended by a variable.

\lean{Copying/PropositionEnvironmentWeakening}{34}{37}{Weakening.privateCell_step}

\lean{Copying/PropositionEnvironmentWeakening}{38}{41}{Weakening.constantResult_step}

\lean{Copying/PropositionEnvironmentWeakening}{42}{47}{Weakening.privateCell_variables}

\lean{Copying/PropositionEnvironmentWeakening}{48}{55}{Weakening.constantResult_variables}

Induction on $h$ now proves
\[
 P_{k,h}\leadsto A_h(V_k(P_{0,0})),\qquad
 Q_{k,h}\leadsto A_h(V_k(Q_{0,0})).
\]
The base is the variable-only result; the successor composes the one-atom law
with the extended induction hypothesis. The named-environment versions take
$k=|\Gamma|$ and $h=|\Delta|$ from the existing environment structure, which
already carries the two distinctness proofs. There is no additional freshness,
termination, machine-determinism or translation-correctness premise.

\lean{Copying/PropositionEnvironmentWeakening}{56}{62}{Weakening.privateCell_environment}

\lean{Copying/PropositionEnvironmentWeakening}{63}{69}{Weakening.constantResult_environment}

\lean{Copying/PropositionEnvironmentWeakening}{70}{74}{Weakening.privateCell_named}

\lean{Copying/PropositionEnvironmentWeakening}{75}{81}{Weakening.constantResult_named}

Each calculation yields exact successful-run equivalence when inserted at
any admitted recursive closing site. The surrounding graph must be focus
free, and filling the source must be hole free; these are the existing
observation admissibility conditions, not unproved machine or template laws.
Apply the calculation's proved site-lifting operation and exact-equivalence
theorem. The general result covers every renamed term; the two program
specializations cover arbitrary $k,h$ at their original depths. Termination
equivalence follows by forgetting the exact count, but no count is discarded
in these declarations.

\lean{Copying/PropositionEnvironmentWeakening}{82}{96}{Weakening.term_weakening_addressed}

\lean{Copying/PropositionEnvironmentWeakening}{97}{105}{Weakening.privateCell_environment_addressed}

\lean{Copying/PropositionEnvironmentWeakening}{106}{112}{Weakening.constantResult_environment_addressed}

\Needspace{12\baselineskip}
\subsection{Checked instances and dependency boundary}

The effectful regression term combines a freely chosen unary and binary
operation with an old variable and dereference of an old atom, in a nonempty
mixed environment. Its theorem specializes full weakening for every operation
and either sort. A second term has two nested abstractions and invokes the
same outer callback twice before reading an old atom; its weakening calculation
checks capture-free lifting under both binders. The self-application example
specializes the theorem to the existing $\Omega$ syntax without a termination
premise; this is not a new proof of divergence. Finally, a product packages
both original program calculations at two variables and three atoms. These
instances check actual translated renamed terms rather than manually
substituted graph endpoints.

\lean{Copying/TermWeakeningExamples}{12}{14}{Weakening.Examples.effectTerm}

\lean{Copying/TermWeakeningExamples}{15}{21}{Weakening.Examples.all_operations}

\lean{Copying/TermWeakeningExamples}{22}{25}{Weakening.Examples.repeatedCallback}

\lean{Copying/TermWeakeningExamples}{26}{32}{Weakening.Examples.nested_repeated_callback}

\lean{Copying/TermWeakeningExamples}{33}{39}{Weakening.Examples.self_application}

\lean{Copying/TermWeakeningExamples}{40}{43}{Weakening.Examples.proposition_mixed_environment}

\paragraph{Dependency map and remaining work.}
Completed dependencies are: injective atomic occurrences $\to$ exhaustive
one-weakening complements $\to$ boundary-fixing atomic isomorphisms; these,
together with the existing unary, boxed and independent padded binary laws,
give the full syntactic induction $\to$ arbitrary finite ambient-extension
calculations for both original programs $\to$ exact addressed observations.
The semantic leaves remain the already proved Idem and BPullW results and
recursive isomorphism laws. Focus-free hosts and hole-free fillings are
explicit application conditions.

Section~\ref{sec:proposition-opening} now proves the actual allocation and outer
beta prefix and transports it through a certified ambient context. The rest
of Figure~22's closed-program equivalence chain is unfinished, as are the
correspondence and further substitution/stability obligations retained in
Appendix~B.1. The ambient-extension theorem does not assume or supply that
closed equivalence, so it cannot by itself prove Proposition~12.2. The estimate
remains 70--80\% of mathematical effort, with low confidence: a substantial
structural task has closed, but the remaining chain and correspondence work
still prevent a defensible narrower estimate.
